\section{Basis / DP}
\subsection{Basis Vector}
\begin{minted}{c++}
template<typename T = int, int LG = 31>
struct Basis {
  T basis[LG]; int sz;
  vector<int> who;
  Basis() {
    fill(basis, basis + LG, 0), sz = 0;
    who.clear();
  }
  void reset() {
    fill(basis, basis + LG, 0), sz = 0;
    who.clear();
  }
  bool insert(T mask, int idx = -1) {
    for (int i = LG - 1; i >= 0; --i) {
      if (mask >> i & 1) {
        if (basis[i]) mask ^= basis[i];
        else {
          basis[i] = mask, ++sz;
          who.push_back(idx);
          return true;
        }
      }
    }
    return false;
  }
  T maxXor() {
    T ans = 0;
    for (int i = LG - 1; i >= 0; --i) {
      ans = max(ans, ans ^ basis[i]);
    }
    return ans;
  }
  bool can(T x) {
    for (int i = LG - 1; i >= 0; --i) {
      if (x >> i & 1) {
        if (!basis[i]) return false;
        x ^= basis[i];
      }
    }
    return true;
  }
  T kth(T k) {
    if (k < 1 || k > ((T)1 << sz)) return -1;
    T mask = 0;
    int64_t tot = 1LL << sz;
    for (int i = LG - 1; i >= 0; --i) {
      if (basis[i]) {
        int64_t low = tot / 2;
        if ((low < k && (~mask >> i & 1)) ||
        (low >= k && (mask >> i & 1))) {
          mask ^= basis[i];
        }
        if (low < k) k -= low;
        tot /= 2;
      }
    }
    return mask;
  }
\end{minted}
\begin{minted}{c++}
  // number of subsets having xor < x
  T count_less(T x) {
    if (x < 0) {
      return 0;
    }
    T ret = 0;
    T cnt = ((T)1 << sz);
    T mask = 0;
    for (int i = LG - 1; i >= 0; i--) {
      if (basis[i]) {
        if (x >> i & 1) {
          ret += (cnt >> 1);
          if (~mask >> i & 1) {
            mask ^= basis[i];
          }
        } else {
          if (mask >> i & 1) {
            mask ^= basis[i];
          }
        }
        cnt >>= 1;
      } else {
        if ((x ^ mask) >> i & 1) {
          if (x >> i & 1) {
            return ret + cnt;
          } else {
            return ret;
          }
        }
      }
    }
    return ret;
  }
  T count_leq(T x) {
    return count_less(x + 1);
  }
  T jumpK(T val, T k) {
    T p = count_leq(val);
    return kth(p + k);
  }
};
\end{minted}
\subsection{Sum of Digits up to N}
\begin{minted}{c++}
uint64_t sumDigitsUpTo(uint64_t n){
  if(n <= 0) return 0;
  uint64_t sum = 0;
  for(uint64_t f = 1; f <= n; f *= 10){
    uint64_t hi = n / (f * 10);
    uint64_t cur = (n / f) % 10;
    uint64_t lo = n % f;
    sum += hi * 45LL * f;
    sum += cur * (lo + 1);
    sum += f * (cur * (cur - 1) / 2);
  }
  return sum;
}
\end{minted}
\subsection{Number of Digits up to N}
\begin{minted}{c++}
int64_t get(int64_t n) {
  int64_t base = 1, digits = 1, now = 0;
  while (base * 10 <= n) {
    base *= 10;
    now += (base - (base / 10 == 0 ? 1 : base / 10))
    * digits;
    digits += 1;
  }
  n -= (base - 1);
  now += n * digits;
  return now;
}
\end{minted}
\subsection{Linear Combination Checking}
Given array \texttt{arr} and integer N, check if N is
a linear combination of arr's elements.

E.g.\ $N = a \cdot x + b \cdot y + c \cdot z + \ldots$

Ans: $N \bmod \text{GCD}(\text{all elements}) == 0$.
\subsection{SOS DP}
\begin{minted}{c++}
void zeta_subset(vector<int64_t> &dp, int bits) {
  for (int b = 0; b < bits; b++) {
    for (int mask = 0; mask < (1 << bits); mask++)
    {
      if (mask >> b & 1) {
        dp[mask] += dp[mask ^ (1 << b)];
      }
    }
  }
}
void zeta_superset(vector<int64_t> &dp, int bits)
{
  for (int b = 0; b < bits; b++) {
    for (int mask = 0; mask < (1 << bits); mask++)
    {
      if (!(mask >> b & 1)) {
        dp[mask] += dp[mask | (1 << b)];
      }
    }
  }
}
\end{minted}
\begin{minted}{c++}
void Solve()
{
  ll n;
  cin >> n;
  vector<ll> arr(n);
  for (auto& x : arr)
    cin >> x;
  ll N = 1 + __lg(*max_element(arr.begin(),
  arr.end()));
  vector<ll> sos2(1 << N, 0), sos(1 << N, 0);
  for (auto it : arr) {
    sos[it]++;
    sos2[((1 << N) - 1) ^ it]++;
  }
  for (int i = 0; i < N; ++i)
    for (int mask = 0; mask < (1 << N); ++mask)
      if (mask & (1 << i))
        sos[mask] += sos[mask ^ (1 << i)],
        sos2[mask] += sos2[mask ^ (1 << i)];
  // the number of elements y such that x|y = x
  // the number of elements y such that x&y = x
  // the number of elements y such that x&y != 0
  for (auto& it : arr) {
    cout << sos[it] << ' ';
    cout << sos2[((1 << N) - 1) ^ it] << ' ';
    cout << n - sos[((1 << N) - 1) ^ it] << ' ';
  }
}
\end{minted}
\subsection{LIS}
\begin{minted}{c++}
int lis(vector<int> const& a) {
  int n = a.size();
  const int INF = 1e9;
  vector<int> d(n+1, INF);
  d[0] = -INF;
  for (int i = 0; i < n; i++) {
    int l = upper_bound(d.begin(), d.end(),
    a[i]) - d.begin();
    if (d[l-1] < a[i] && a[i] < d[l])
      d[l] = a[i];
  }
  int ans = 0;
  for (int l = 0; l <= n; l++) {
    if (d[l] < INF)
      ans = l;
  }
  return ans;
}
\end{minted}
\subsection{LIS up to Every Index}
\begin{minted}{c++}
vector<int>
longestObstacleCourseAtEachPosition(vector<int
>& a) {
  vector<int> lis,ans;
  for(auto &it:a){
    int indx =
    upper_bound(lis.begin(),lis.end(),it)-lis.begin();
    if(indx==(int)lis.size()) lis.push_back(it);
    else lis[indx] = it;
    ans.push_back(indx+1);
  }
  return ans;
}
\end{minted}
\subsection{Subset Sum Bitset}
// Only works when the sum of values is not
// greater than N.
// TC: O(N / 32 * sqrt(N))
\begin{minted}{c++}
const int N = 1e6 + 6;
bitset<N> dp;
dp.set(0);
for (auto &[w, c] : comp) { // w is the value and c
// is the count of it's
  for (int i = 0; 1 << i <= c; ++i) {
    dp |= dp << (w * (1 << i));
    c -= 1 << i;
  }
  dp |= dp << (w * c);
}
dp.test(k) // if it's possible to make sum equal to k
\end{minted}

OR convolution can be done using SOS DP.
\subsection{All Submasks in \texorpdfstring{$O(3^N)$}{O(3^N)}}
\begin{minted}{c++}
for (int mask = 0; mask < 1 << n; mask++) {
  for (int sub = mask; sub > 0;
  sub = (sub - 1) & mask) {
  }
}
\end{minted}
\subsection{Max Rectangular Sum}
\begin{minted}{c++}
int main() {
  int n;
  while (scanf("\%d", &n) != EOF) {
    int mx = -1280, arr[n + 1][n + 1], linear[n],
    gsum, csum, qsum[n][n];
    for (int i = 0; i < n; i++) {
      for (int j = 0; j < n; j++) {
        scanf("\%d", &arr[i][j]);
        if (j == 0)
          qsum[i][j] = arr[i][j];
        else
          qsum[i][j] = arr[i][j] + qsum[i][j - 1];
      }
    }
    gsum = arr[0][0];
    for (int j = 0; j < n; j++) {
      for (int k = j; k < n; k++) {
        for (int i = 0; i < n; i++) {
          int temp = qsum[i][k] - qsum[i][j] +
          arr[i][j];
          if (i == 0) {
            csum = temp;
            continue;
          }
          csum = max(temp, csum + temp);
          gsum = max(csum, gsum);
        }
      }
    }
    printf("\%d\n", gsum);
  }
  return 0;
}
\end{minted}
\subsection{Derangements}
\begin{minted}{c++}
if (n > 0) D[1] = 0;
if (n > 1) D[2] = 1;
for (int i = 3; i <= n; i++) {
  D[i] = (i - 1) * (D[i - 1] + D[i - 2]);
}
\end{minted}
\subsection{Stars and Bars}
 $C(n + k - 1, k)$;
\subsection{Catalan Numbers}
$n=$ number of pairs: $C(2n, n) - C(2n, n-1)$;

$n=$ number of brackets: $\frac{1}{n+1} \cdot C(2n, n)$
\subsection{Lucas Theorem}
\begin{minted}{c++}
long long lucas(long long n, long long r) {
  if (r == 0) return 1;
  return lucas(n / MOD, r / MOD) * nCr(n %
  MOD, r % MOD) % MOD;
}
\end{minted}
\subsection{Binomial Coefficient Properties}
\[
\binom{n}{k} = \binom{n}{n-k}
\]
\[
\sum_{k=0}^{m} \binom{n-1}{k-1} = \binom{n+m+1}{m}
\]
\[
\binom{n}{k} = \binom{n-1}{k-1} + \binom{n-1}{k}
\]
\[
\sum_{k=0}^{n} \binom{n}{k} = 2^n
\]
\[
\binom{n}{0} + \binom{n}{1} + \ldots + \binom{n}{n} = 2^n
\]
\[
\binom{n}{0} + \binom{n}{2} + \binom{n}{4} + \ldots = 2^{n-1}
\]
\[
1\binom{n}{1} + 2\binom{n}{2} + \ldots + n\binom{n}{n} = n \cdot 2^{n-1}
\]

Connection With Fibonacci numbers:
\[
\binom{n}{0} + \binom{n-1}{1} + \ldots + \binom{n-k}{k} + \ldots + \binom{0}{n} = F_{n+1}
\]

